% Complete Section 8.7, original PDF 114-117 / printed 104-107.
% block: stokes-equations-heading page:114 type:section
\section{{Stokes 方程}}\label{sec:stokes-equations}

% block: neumann-mixed-formulation pages:114,115 type:ordinary-paragraph
{回顾 $\mathbb R^n$ 中区域 $\Omega$ 上 Poisson 方程的 Neumann 问题的混合弱形式. 给定 $f\in\hat L^2$, 求 $u\in\mathring H(\operatorname{div},\Omega)$, $p\in\hat L^2(\Omega)$, 使得}
\begin{equation}\label{eq:neumann-poisson-mixed-formulation}
\begin{aligned}\langle u,v\rangle-\langle p,\operatorname{div}v\rangle&=0,&v&\in\mathring H(\operatorname{div}),\\\langle\operatorname{div}u,q\rangle&=\langle f,q\rangle,&q&\in\hat L^2.\end{aligned}
\end{equation}
{($\hat L^2$ 中的帽号表示与常数正交的子空间, 反映了 $f$ 的积分必须为零, 且 $p$ 只确定到相差一个加性常数. 也可以在 $L^2$ 中工作, 并为调和形式空间, 即常数, 增加相应项.) 问题 \eqref{eq:neumann-poisson-mixed-formulation} 是下述片段关联的 Hodge Laplacian 的混合弱形式}
\begin{equation}\label{eq:neumann-divergence-complex}
\mathring H(\operatorname{div})\xrightarrow{\operatorname{div}}\hat L^2(\Omega)\to0,
\end{equation}
{它取自带边界条件的 $L^2$ de Rham 复形末端. 利用 Chapter~\ref{chap:finite-element-differential-forms} 发展的有限元微分形式, 很容易离散化此形式. 下面将此问题与 Stokes 流问题比较.}

% block: stokes-strong-problem-boundaries page:115 type:ordinary-paragraph
{Stokes 流方程的强形式, 是求表示速度的向量场 $u$ 和表示压力的标量场 $p$, 满足}
\[
-\mu\Delta u+\operatorname{grad}p=f,\qquad\operatorname{div}u=0\quad\text{in }\Omega.
\]
{这里 $\mu$ 是动力黏度, 向量场 $f$ 给出所施加的体力密度. 一个简单的边界条件选择是无滑移条件, 即 $\partial\Omega$ 上 $u=0$. 更现实地, 可以在部分边界上取无滑移条件, 在其余边界上取应力边界条件 $2\mu\epsilon(u)n-pn=0$, 该条件用 $u$ 的对称梯度 $\epsilon(u)$ (应变率) 和外单位法向量 $n$ 陈述.}

% block: stokes-weak-formulation page:115 type:ordinary-paragraph
{为简单起见, 仍采用纯无滑移边界条件, 弱形式是求 $u\in\mathring H^1(\Omega;\mathbb R^n)$, $p\in\hat L^2(\Omega)$, 使得}
\begin{equation}\label{eq:stokes-weak-formulation}
\begin{aligned}b(u,v)-\langle p,\operatorname{div}v\rangle&=\langle f,v\rangle,&v&\in\mathring H^1(\Omega;\mathbb R^n),\\\langle\operatorname{div}u,q\rangle&=0,&q&\in\hat L^2(\Omega).\end{aligned}
\end{equation}
{\eqref{eq:stokes-weak-formulation} 中的双线性形式 $b:\mathring H^1(\Omega;\mathbb R^n)\times\mathring H^1(\Omega;\mathbb R^n)\to\mathbb R$ 由下式给出}
\begin{equation}\label{eq:stokes-strain-bilinear-form}
b(u,v)=2\mu\langle\epsilon(u),\epsilon(v)\rangle.
\end{equation}
{利用无滑移边界条件和无散条件 $\operatorname{div}u=0$, 可将其化为 $b(u,v)=\mu\langle\operatorname{grad}u,\operatorname{grad}v\rangle$, 但我们不愿这样做, 因为在部分边界上施加应力 (或其他) 边界条件时, 就不再能这样化简.}

% block: stokes-smoothed-complex pages:115,116 type:ordinary-paragraph
{形式 \eqref{eq:stokes-weak-formulation} 与 Neumann 问题的混合弱形式 \eqref{eq:neumann-poisson-mixed-formulation} 共有项 $\langle p,\operatorname{div}v\rangle$ 和 $\langle\operatorname{div}u,q\rangle$. 然而, \eqref{eq:stokes-weak-formulation} 在 $\mathring H^1(\Omega;\mathbb R^n)$ 中求 $u,v$, 而 \eqref{eq:neumann-poisson-mixed-formulation} 则在 $\mathring H(\operatorname{div})$ 中求它们. 这启发了一个稍有不同的复形:}
\begin{equation}\label{eq:stokes-divergence-complex}
\mathring H^1(\Omega;\mathbb R^n)\xrightarrow{\operatorname{div}}\hat L^2(\Omega)\to0.
\end{equation}
{复形 \eqref{eq:stokes-divergence-complex} 与 \eqref{eq:neumann-divergence-complex} 的不同在于, 对复形的第一个空间要求更高的光滑性. 为完整定义该复形, 除 \eqref{eq:stokes-divergence-complex} 所示的定义域空间外, 还必须指定 $\operatorname{div}$ 作用的底层空间. 事实上, 将 \eqref{eq:stokes-divergence-complex} 视为\emph{有界} Hilbert 复形, 因而底层空间与定义域空间重合. 在第一个空间 $\mathring H^1(\Omega;\mathbb R^n)$ 上, 以双线性形式 $b$ 为内积, 它与该空间上的 $H^1$ 内积等价. 作出这些选择后, 弱形式 \eqref{eq:stokes-weak-formulation} 就是复形 \eqref{eq:stokes-divergence-complex} 的抽象 Hodge Laplacian 的混合弱形式. 因此, 只要 Hilbert 复形是闭的且没有调和形式, 此问题的适定性就来自 Chapter~\ref{chap:hilbert-complexes} 的一般理论. 在当前情形中, 这意味着要求散度算子将 $\mathring H^1(\Omega;\mathbb R^n)$ 满射到 $\hat L^2(\Omega)$. 这比 $\operatorname{div}$ 将 $\mathring H(\operatorname{div},\Omega)$ 满射到 $\hat L^2(\Omega)$ 更强, 后一结论在 Section~\ref{sec:vector-operators-adjoints} 中已容易地建立. 这一更强结果广泛用于 Stokes 方程的分析, 已有多种证明, 适用于一般性程度不同的区域. 关于不同方法和结果的出色讨论, 见 [\BookNumber{35}, Section 2.2].}

% block: stokes-discrete-subcomplex-projections page:116 type:ordinary-paragraph
{应用 Chapter~\ref{chap:approximation-hilbert-complexes} 的理论, 得到准确地对 \eqref{eq:stokes-divergence-complex} 作 Galerkin 离散化的条件. 只需 Galerkin 子空间 $V_h\subset\mathring H^1(\Omega;\mathbb R^n)$ 和 $S_h\subset L^2(\Omega)$ 满足子复形性质}
\begin{equation}\label{eq:stokes-discrete-subcomplex}
\operatorname{div}V_h\subset S_h
\end{equation}
{并且存在有界上链投影:}
\[
\pi_V:\mathring H^1(\Omega;\mathbb R^n)\to V_h,\quad\pi_S:L^2(\Omega)\to S_h,\quad\operatorname{div}\pi_Vv=\pi_S\operatorname{div}v.
\]
{于是由 FEEC 理论或更初等的论证可得, Galerkin 离散化稳定且满足误差估计}
\begin{equation}\label{eq:stokes-galerkin-error}
\|u-u_h\|_1+\|p-p_h\|_0\leq c\left(\inf_{v\in V_h}\|u-v\|_1+\inf_{p\in S_h}\|p-q\|_0\right).
\end{equation}

% block: stokes-exact-mass-conservation page:116 type:ordinary-paragraph
{满足子复形性质 $\operatorname{div}V_h\subset S_h$ 的 Stokes 单元对, 还有一个重要性质. 离散速度 $u_h$ 守恒质量, 即它精确求解连续性方程 $\operatorname{div}u_h=0$. 这立即由 Galerkin 方程}
\[
\langle\operatorname{div}u_h,q\rangle=0
\]
{中取 $q=\operatorname{div}u_h\in S_h$ 得到. 质量守恒对许多应用都是非常理想的性质, 尤其在处理非线性和/或含时 Navier–Stokes 方程时. 近期综述 \BookCite{55} 对此作了深入讨论. 经典 Stokes 单元对 (MINI 单元, $\mathcal P_2$--$\mathcal P_0$, Taylor–Hood 等), 例如 \BookCite{45} 或 [\BookNumber{19}, Chapter 8] 讨论的单元对, \emph{不}满足子复形性质. 虽然它们满足估计 \eqref{eq:stokes-galerkin-error}, 却不守恒质量.}

% block: two-dimensional-stokes-complex pages:116,117 type:ordinary-paragraph
{近十年左右, 质量守恒单元的发展一直是活跃的研究领域. 主要方法是发展片段 \eqref{eq:stokes-divergence-complex} 的有限元子复形, 以及以它为一部分的更长复形的有限元子复形. 在二维中, 所讨论的完整复形为}
\[
0\to H^2(\Omega)\xrightarrow{\operatorname{curl}}H^1(\Omega;\mathbb R^n)\xrightarrow{\operatorname{div}}L^2(\Omega)\to0,
\]
{这里为简化起见, 我没有施加边界条件, 也没有去除调和形式. 这个 \emph{Stokes 复形}是 $L^2$ de Rham 复形的光滑化版本. 虽然 Stokes 方程的离散化只需要复形中最后两个非零空间的子空间, 但在构造和分析这些子空间时, 考察整个复形及其子复形往往有所帮助. 这就将 Stokes 方程的质量守恒有限元与 $H^2(\Omega)$ 的有限元子空间联系起来, 后者是有限元领域长期研究的一个有挑战性的问题. Mardal, Tai 和 Winther 于 2002 年给出了二维 Stokes 复形的首个有限元离散化 \BookCite{57}. 然而它是非协调离散化, 因而不是真正的子复形. 2013 年, Guzmán 和 Neilan 构造了在一般网格上的首个协调有限元离散化 \BookCite{49}. 他们使用由某些有理函数增广的低次分片多项式. 同年 Falk 和 Neilan \BookCite{42} 构造了一族 Stokes 复形的协调分片多项式子复形. 此外还有许多其他可能的单元选择. 更多信息见 [\BookNumber{55}, Figure 4.1].}

% block: three-dimensional-stokes-complex page:117 type:ordinary-paragraph
{三维中质量守恒 Stokes 单元的发展, 也可类似地基于\emph{三维 Stokes 复形}}
\[
0\to H^2\xrightarrow{\operatorname{grad}}H^1(\operatorname{curl})\xrightarrow{\operatorname{curl}}H^1(\Omega;\mathbb R^3)\xrightarrow{\operatorname{div}}L^2\to0,
\]
{其中 $H^1(\operatorname{curl})$ 由所有其旋度仍属于 $H^1$ 的 $H^1$ 向量场组成. Neilan 在 \BookCite{60} 中构造了一个有限元子复形. 最近 Guzmán, Fu 和 Neilan \BookCite{48} 针对特殊的网格几何发展了另一种方法.}
